\chapter{2002 Nobel Prize in Economics}
\textbf{Laureates: }Daniel Kahneman (Israel/USA), Vernon Smith (USA)

\section*{Reason for Award}
For having integrated insights from psychological research into economic science

\section{What They Did}
Kahneman and Smith transformed economics by incorporating psychological
realism into economic analysis, challenging the standard model of the rational
economic agent that had dominated the discipline for decades.

Kahneman, together with the late Amos Tversky, developed prospect theory,
which describes how people actually make decisions under risk and uncertainty.
Unlike expected utility theory, which assumes that individuals evaluate
outcomes based on final wealth positions and maximize expected utility,
prospect theory shows that people evaluate outcomes relative to a reference
point--typically the status quo. Losses relative to the reference point are
weighted more heavily than equivalent gains, a phenomenon known as loss
aversion. Empirical estimates suggest that losses feel roughly twice as
painful as equivalent gains feel pleasurable.

Kahneman and Tversky identified a range of cognitive heuristics and biases
that affect judgment under uncertainty. The availability heuristic leads
people to judge the probability of events based on how easily examples come
to mind. The representativeness heuristic leads people to judge probabilities
based on similarity to stereotypes, ignoring base rates and sample sizes. The
anchoring heuristic causes people to rely too heavily on initial information
when making decisions.

Smith established experimental economics as a legitimate scientific method.
Prior to his work, economists largely treated laboratory experiments as
irrelevant to understanding real-world economic behavior. Smith demonstrated
that markets can achieve efficient outcomes even when participants have
limited information and face various frictions. In his pioneering double-auction
experiments, subjects traded a perishable asset over multiple periods, and
prices converged rapidly to competitive equilibrium values. This validated
the invisible hand hypothesis while also revealing systematic deviations that
informed behavioral theory.

\section{Impact on Economic Thinking}
Kahneman's behavioral revolution challenged the rational agent paradigm that
had dominated economics since the neoclassical revolution. Prospect theory
provided a descriptive alternative to expected utility theory that could
explain phenomena such as the equity premium puzzle, the disposition effect
in trading, underinvestment in retirement savings, and the endowment effect.

Smith's experimental methods transformed economics from a discipline that
relied primarily on theoretical deduction and observational data into one
that could test hypotheses under controlled conditions. Experimental economics
became a recognized subfield with its own journals, professional associations,
and methodological standards.

The synthesis of behavioral insights and experimental methods created new
interdisciplinary research programs. Behavioral game theory incorporated
psychological realism into strategic analysis, relaxing the assumption of
common knowledge of rationality. Neuroeconomics investigated the neural
correlates of decision-making using tools like functional magnetic resonance
imaging (fMRI) and electroencephalography (EEG).

Kahneman's distinction between the experiencing self and the decision-making
self, developed later in his career, raised profound questions about welfare
economics and the measurement of well-being that challenged traditional
cost-benefit analysis. His work also influenced the design of choice
architectures in organizations and government agencies.

\section{Modern Perspectives Today}
Kahneman and Smith's legacy is evident throughout contemporary economics and
public policy. Nudge theory, popularized by Thaler and Sunstein, applies
behavioral insights to policy design through choice architecture--altering the
way choices are presented to influence behavior without restricting freedom of
choice. Governments worldwide have established behavioral insight units to
improve policy outcomes in areas ranging from tax collection and organ donation
to retirement savings and energy conservation.

Smith's experimental traditions have expanded beyond laboratory settings to
field experiments and natural experiments that test interventions in real-world
contexts. Randomized controlled trials (RCTs) now inform development economics,
education policy, healthcare delivery, and environmental policy.

Behavioral economics has influenced financial regulation, consumer protection,
and retirement savings policy. The designation of "behavioral economics" as a
major research area by the American Economic Association reflects the field's
institutional recognition. Their combined contributions remain central to
understanding how people actually make economic decisions.
